% !TEX root = ../main.tex
\ulfgchapter{1}{General Introduction}{chp:intro}

\section{Motivation}

Automotive radar has a distinctive role in the sensor suite of a modern vehicle. Cameras can degrade
under poor illumination, glare and reduced visibility. Lidar can be affected by rain, snow and
airborne particles. Millimetre-wave radar is comparatively robust under many of these conditions.
It also measures radial velocity directly, from the Doppler shift, rather than estimating it by
comparing successive frames. For a driver assistance system that must distinguish an object closing
at \SI{2}{\metre\per\second} from one closing at \SI{20}{\metre\per\second}, a direct velocity
measurement is valuable.

The frequency-modulated continuous-wave (FMCW) waveform makes this practical at low cost. The radar
transmits a linear frequency ramp, called a chirp, and mixes the received echo with the transmitted
signal. The result is a low-frequency beat tone. Its frequency gives the target range, and its phase
change from one chirp to the next gives the radial velocity. Two Fourier transforms turn a block of
received samples into a \emph{range--Doppler} (RD) map: a two-dimensional image in which a target
appears as a concentration of energy at its range and its radial velocity. Everything downstream ---
clustering, tracking, classification, and ultimately the driving decision --- works from what is
extracted from this map.

Extracting targets from the RD map is a detection problem, and it is the first stage of that chain.
Detection determines which measurements are passed to later processing. A tracker can bridge an
occasional missed detection, but a target that is missed persistently cannot be tracked, classified
or acted upon, because no measurement of it reaches those stages. Excessive false detections cause
the opposite problem: they consume tracking resources and can trigger unnecessary responses.
Detection therefore places a strong limit on the performance of everything built above it.

For about fifty years the standard solution has been the constant false alarm rate (CFAR) family of
detectors. The idea is simple and remains effective. Comparing each cell against a fixed threshold
would require knowing the noise power, which changes with temperature, terrain, weather and
interference. Instead, CFAR estimates the local background from cells surrounding the cell under
test (CUT), and scales that estimate by a multiplier chosen to hold the false-alarm rate constant.
The detector therefore adapts to conditions it was never told about. CFAR is analytically tractable,
needs no training data, is cheap to compute, and its behaviour can be predicted in closed form under
stated assumptions.

Those assumptions are also its limitation, and this was characterised early. Gandhi and
Kassam~\cite{gandhi1988cfar} showed that cell-averaging CFAR degrades in situations that matter in
practice. At a clutter edge, the reference window covers two different background levels and the
threshold is set from the wrong one. Under target masking, a second target inside the reference
window raises the background estimate and suppresses detection of the first. The main variants ---
greatest-of, smallest-of and order-statistic --- each improve one failure mode and worsen another.
Their conclusion was that no single CFAR detector performs well across all backgrounds, and that an
adaptive approach may be needed when operating conditions change.

That conclusion invites learning the detector from data, and a substantial body of recent work has
done exactly that. Convolutional networks applied to RD maps have been reported to outperform CFAR,
with improvements ranging from a fraction of a decibel to large reductions in false-alarm rate.

\section{Problem statement}
\label{sec:problem}

These reported improvements are difficult to compare with one another, and difficult to interpret
individually.

A comparison between a classical and a learned detector is informative only if both receive
comparable information, are evaluated at the same operating point, and are scored by the same rule.
Across the reviewed literature, these conditions are not always simultaneously satisfied. Five
differences recur:

\begin{itemize}
\itemsep2pt
\item \textbf{Input preprocessing.} The network may receive a windowed, matched-filtered or
multi-resolution input, while the CFAR baseline receives raw magnitude data.
\item \textbf{Spatial context.} The network may observe a much larger region than the CFAR reference
window, so the two are not solving the same local problem.
\item \textbf{Retraining.} The network may be retrained for each test condition, while the classical
detector keeps a single global parameter set.
\item \textbf{Operating point.} False-alarm rates are often matched by \emph{design}, using the
analytic CFAR expression, rather than by measurement. CFAR detectors do not reach their design
$\Pfa$ in clutter, a point reported by authors who also report learned
wins~\cite{yavuz2021radcnn}.
\item \textbf{Scoring tolerance.} A detection may count as correct if it falls within some distance
of the true position. A detector whose response is spread over a broad region then benefits from
that tolerance.
\end{itemize}

Each of these choices can be reasonable on its own. Taken together, they make the reported margin
depend partly on the experimental setup rather than only on the detector. Three questions that
should be separable then become hard to answer:

\begin{enumerate}
\itemsep2pt
\item \textbf{Is there a genuine advantage} when both detectors receive comparable information and
are compared at a matched operating point?
\item \textbf{Where does it come from?} A margin that is not attributed to a mechanism is not
understood, and there is no basis for expecting it to reappear under different conditions.
\item \textbf{Does it survive outside the training domain?} A detector trained on simulated data and
evaluated on the same simulator has shown that it learned the simulator. Whether it learned
\emph{radar} is a separate question, and one that much of the literature does not test.
\end{enumerate}

The third question is especially relevant for a learned detector, for a reason specific to this
problem. CFAR estimates its threshold from the data in front of it at inference time, so its
adaptation mechanism travels with it into any domain. A conventionally frozen feed-forward network
has no explicit mechanism to re-estimate its decision rule when the deployment distribution changes;
its behaviour is fixed by weights fitted to the training distribution.

This gives the project its central framing. The purpose is not only to determine whether a small
convolutional detector can outperform CFAR, but to determine what a small learned local decision
rule adds when both detectors operate over comparable range--Doppler context.

\section{Objectives}

Four objectives follow from the three questions above.

\begin{description}
\itemsep4pt
\item[O1 --- Build a controlled comparison.] Construct a classical arm and a learned arm that
receive the same input representation, observe comparable local context, and are scored by a single
evaluator at a matched \emph{measured} false-alarm rate. The learned detector's receptive field is
deliberately matched to the spatial footprint of the CFAR window, providing comparable local context
while allowing each detector to process that context according to its own decision rule. The two
rules are not identical: CFAR uses the CUT as its test statistic, excludes the surrounding guard
cells, and estimates the interference level from the remaining reference cells, whereas the network
may use all cells in its footprint jointly.

\item[O2 --- Measure the margin.] Quantify the difference across a set of conditions chosen to
include both regimes where CFAR is close to optimal and regimes where it is known to fail. Report
confidence intervals, and correct for multiple comparisons.

\item[O3 --- Attribute the margin.] Determine which mechanisms account for any margin by testing
candidate explanations experimentally rather than arguing for them. Two explanations in the
literature both predict a learned advantage in homogeneous conditions, so neither can be accepted
without testing the other. These experiments are designed to support or rule out specific
explanations, and so narrow the set of plausible mechanisms; they are not expected to establish a
single cause conclusively.

\item[O4 --- Test external validity.] Evaluate both frozen detectors, with no retraining or
recalibration of any kind, on real automotive radar data taken from a public dataset, under a
protocol registered and hash-frozen before the test data were scored.
\end{description}

One objective is deliberately excluded. This project does not aim to produce a state-of-the-art
detector. A larger network, trained on real data with modern augmentation, would very likely detect
more targets than the one built here. The aim is to measure a difference under conditions where the
measurement can be trusted, and for that purpose a small, well-understood result is more useful than
a large unexplained one.

\section{Contributions}

\begin{enumerate}
\itemsep5pt
\item \textbf{A controlled like-for-like protocol.} The learned detector's receptive field is
matched to the spatial footprint of the CFAR window; both arms are scored by one evaluator through a
single code path; operating points are matched at measured rather than design $\Pfa$; and results
are reported across a tolerance ladder rather than at a single gate. Each choice addresses one of
the asymmetries listed in Section~\ref{sec:problem}, and each is documented together with the
asymmetry it addresses.

\item \textbf{A quantified and externally calibrated margin.} On synthetic data, the learned
detector reaches a detection probability of 50\,\% at a cell signal-to-noise ratio
\SynSnrGainDL{}\,dB lower than CA-CFAR, at matched measured $\Pfa$. In other words, it can tolerate
a weaker target signal while maintaining the same detection level. At clutter edges the advantage
reaches \SynClutterEdgeDiff{} in $\Pd$, which falls inside the range reported independently for
clutter transitions on real X-band measurements~\cite{gusland2020deeptemporal}. The CFAR loss
proposed as an explanation for learned wins in homogeneous conditions~\cite{lallouet2023loss} is
measured here rather than assumed, at \SynCfarLossCA{}\,dB relative to a clairvoyant bound.

\item \textbf{An attribution programme} that decomposes the margin into named candidate mechanisms,
each paired with an experiment capable of falsifying it, rather than reporting the margin as a
single unexplained number.

\item \textbf{A pre-registered zero-shot external validation.} The frozen detectors are evaluated on
the CARRADA automotive radar dataset~\cite{ouaknine2021carrada} with no real-data training,
fine-tuning or recalibration, under a protocol fixed by cryptographic hash before any test frame was
scored. None of the directly comparable studies identified in the reviewed literature performs the
same zero-shot synthetic-to-real evaluation. \textbf{The ordering established on synthetic data
reverses on real data}, and this report treats that reversal as a principal result rather than as an
anomaly to be explained away.

\item \textbf{A reproducibility framework supporting these claims:} frozen model and detector
artefacts, isolation manifests that protect completed experiments from later modification, immutable
raw evidence kept alongside the derived tables, and a corrections log recording every analysis error
found and superseded during the work.
\end{enumerate}

\section{Report outline}

The report follows the order in which the work was carried out, because several conclusions depend
on what was known at the time a decision was taken.

\textbf{Chapter~\ref{chp:background}, \emph{Radar Background and State of the Art}}, develops the
material needed for the rest of the report: the FMCW waveform, range and Doppler estimation, the
construction of the RD map, and the CFAR family with its documented failure modes. It then reviews
learned detection and argues that the published literature separates into families whose metrics are
not interchangeable --- a distinction the remainder of the report follows strictly.

\textbf{Chapter~\ref{chp:simulator}, \emph{Simulator, Radar Specification and Dataset Design}},
describes the simulator, the radar parameters and the labelled datasets. The radar parameters were
chosen to reproduce the RD bin geometry of the real dataset used later, so that a target occupies a
comparable number of cells in both domains. Without this, a transfer failure could be dismissed as a
scaling artefact.

\textbf{Chapter~\ref{chp:detectors}, \emph{Detectors and Evaluation Protocol}}, specifies both
detectors and fixes the evaluation protocol: peak extraction, matching, the operating-point
definition and the statistical treatment. The protocol is frozen before any comparative result is
produced.

\textbf{Chapter~\ref{chp:synthetic}, \emph{Synthetic Benchmark Results}}, reports the benchmark
across the seven conditions, including detection probability as a function of cell signal-to-noise
ratio, which expresses the result in the units used by the comparable literature.

\textbf{Chapter~\ref{chp:attribution}, \emph{Attribution}}, asks where the measured margin comes
from and reports the experiments that separate the competing explanations.

\textbf{Chapter~\ref{chp:realdata}, \emph{Real-Data External Validation}}, evaluates the frozen
detectors zero-shot on real automotive radar data under the pre-registered protocol. It audits the
difference between the two domains before any scoring, reports the reversal, and places both halves
of the project against the published literature.

The \textbf{General Conclusion} states what the evidence supports, what it does not, and which
continuations follow naturally. Three appendices record the corrections log, the frozen
pre-registration and the supporting tables.
